\section{Frobenius comparison and smooth nonvanishing}
\label{proj:fr:section}

We now turn the two opposite jet estimates into a contradiction.
The reduction to positive characteristic is used only for this comparison:
no minimal-model statement or pseudo-effectivity assertion will be
specialized to positive characteristic.

\begin{theorem}[The smooth nonvanishing step]
\label{proj:thm:smooth-nonvanishing}
Assume Assumption~\ref{proj:asm:iitaka} and the lower-dimensional
real-boundary good-model hypothesis of Assumption~\ref{proj:mmp:lower-models}.
If \(X\) is a smooth projective complex variety of dimension \(n\) and
\(K_X\) is pseudo-effective, then \(\kappa(X,K_X)\geq0\).
\end{theorem}

The assertion is immediate in dimension zero. On a smooth curve,
pseudo-effectivity of \(K_X\) gives \(g(X)\geq1\), hence
\(h^0(X,K_X)=g(X)>0\). We therefore suppose \(n\geq2\) and argue by
contradiction. We use the late terminal model \(Y\), the divisors
\(K=K_Y\) and \(A=A_Y\), and the scaling models of the preceding
section. Fix \(m>0\) with \(mK\) Cartier, and write
\[
 \mu_t=\vol(K_X+tA_X)^{1/n},\qquad
 L=r(K+tA).
\]
As before, \(r\) is a positive integer chosen so that
\(r\mu_t\longrightarrow\epsilon\) as \(t\downarrow0\).
Choose the rational number \(\epsilon>0\) sufficiently small for
Propositions~\ref{proj:prop:scalar-jets} and~\ref{proj:prop:large-seshadri}, and also
so that
\begin{equation}
\label{proj:fr:epsilon}
 (14\epsilon)^n<\frac14,\qquad
 80\epsilon<\frac34.
\end{equation}
The inequalities are strict. We next fix a sufficiently large integer
\(q\) as in Proposition~\ref{proj:prop:large-seshadri}, and then fix a
sufficiently small positive rational \(t\). In particular, we may require
\(r>qm+1\). All subsequent characteristic-zero choices will be made with
these data fixed.

\subsection{A small polarization and an actual moving curve}

Choose a smooth common resolution \(W\) of \(Y\) and the scaling model
for \(K+tA\). In this section \(K,A,L\) also denote their pullbacks to
\(W\), and
\[
 K_W=K+E,\qquad E\geq0.
\]
Let \(H_0\) be the pullback of the nef semiample positive part \(N_0\) of
\(L\) on its scaling model. The comparison of canonical pullbacks and
moving parts gives
\begin{equation}
\label{proj:fr:limit-intersections}
 H_0^n\longrightarrow\epsilon^n,\qquad
 LH_0^{n-1}=H_0^n,\qquad
 K_WH_0^{n-1}=KH_0^{n-1}\leq\frac{H_0^n}{r}.
\end{equation}
The maps \(Y\dashrightarrow Y_t\) are small by the choice of the late
model, so every divisor exceptional over \(Y\) is also exceptional
over \(Y_t\).
Indeed, the differences being discarded are exceptional over the
scaling model, and the transform of \(A\) is effective up to rational
linear equivalence. Their intersections with \(H_0^{n-1}\) therefore
give the displayed equalities and inequality. For a fixed ample divisor
\(H_{\mathrm{amp}}\), choose a sufficiently small positive rational
\(\eta\) and put \(H=H_0+\eta H_{\mathrm{amp}}\). By continuity, after
the preceding choices of \(t\) and \(r\), we have
\begin{equation}
\label{proj:fr:intersection-bounds}
 H^n\leq(2\epsilon)^n,\qquad
 K_WH^{n-1}\leq\frac{2H^n}{r},\qquad
 (2L+qmK)H^{n-1}\leq4H^n.
\end{equation}

The canonical divisor of \(W\) is pseudo-effective. Generic
semipositivity, applied with empty boundary
\cite[Theorem~2.1]{Proj-CampanaPaun2015}, and restriction to a sufficiently
general complete-intersection curve give the following fixed data.
Choose \(l>0\) such that \(h=lH\) is integral and sufficiently very
ample, and choose a smooth complete-intersection flag
\[
 W=W_n\supset W_{n-1}\supset\cdots\supset W_1=C,
 \qquad W_{i-1}\in |h|_{W_i},
\]
for which
\begin{equation}
\label{proj:fr:generic-nef}
 \mu_{\min}(\Omega_W^1|_C)\geq0.
\end{equation}
Here and below slopes on \(C\) mean degree divided by rank. One may
obtain the flag by applying restriction to the finitely many
Harder--Narasimhan quotients of \(\Omega_W^1\).

Choose also a very general point \(x\in W\), away from the exceptional
loci, and let \(\pi_x:\widehat W\to W\) be its blowup, with exceptional
divisor \(J\). Set
\[
 D^+=2r(K+2tA)+2E.
\]
This divisor is big. Its sections, pushed to \(Y\), are the scalar
sections in Proposition~\ref{proj:prop:scalar-jets}; the added exceptional
part does not change them. If
\(\rho=\vol(2r(K+2tA))^{1/n}\), then
\[
 \rho\leq4r\mu_t\leq8\epsilon
\]
for our sufficiently small \(t\). Thus every section of a sufficiently
divisible multiple of \(D^+\) has normalized order at \(x\) at most
\(4\rho\leq32\epsilon\).

It follows that \(\pi_x^*D^+-40\epsilon J\) is not pseudo-effective.
Otherwise its convex combination with the big class \(\pi_x^*D^+\)
would make \(\pi_x^*D^+-cJ\) big for some rational
\(32\epsilon<c<40\epsilon\), giving a section of excessive order.
Movable-curve duality \cite[Theorem~2.2]{BDPP} therefore supplies an
actual covering curve class \(\gamma\) on \(\widehat W\) such that
\begin{equation}
\label{proj:fr:curve-test}
 (\pi_x^*D^+)\cdot\gamma
       <40\epsilon\,J\cdot\gamma,\qquad J\cdot\gamma>0.
\end{equation}
For precision, the strict negative pairing can first be detected by a
strongly movable curve: take the pushforward of a general complete
intersection of very ample divisors on one fixed smooth birational
model of \(\widehat W\). Fix this model, its divisors, and the resulting
covering family. This choice is finite algebraic data, rather than a
limiting real movable class. Positivity of \(J\cdot\gamma\) follows
from the strict inequality and pseudo-effectivity of \(\pi_x^*D^+\).

\subsection{Fixed jets and reduction modulo primes}

On \(W\times W\), let
\[
 Z=\PP(mK_1\oplus mK_2),\qquad
 \xi=\OO_Z(1),\qquad
 M=L_1+L_2+q\xi,
\]
with the symmetric-power convention for the projective bundle.
Proposition~\ref{proj:prop:large-seshadri} gives a semiample big model whose
Seshadri constant exceeds \(2n+2\) at a very general torus point.
On the open set where its birational contraction is an isomorphism,
the jet interpretation of the Seshadri constant consequently gives
an integer \(k_0>0\) such that \(|k_0M|\) generates jets of order
at least \((2n+2)k_0\) at a fixed smooth torus point \(z_*\in Z\).
Indeed, on the ample model choose a rational number \(c\) strictly
between \(2n+2\) and its Seshadri constant. On the blowup of the
selected smooth point, the pullback of the ample class minus \(c\)
times the exceptional divisor is ample. Serre vanishing and the
exceptional-divisor sequence then give jets of order \(kc-1\) for
all sufficiently divisible large \(k\). Pulling sections back on the
isomorphic open gives the asserted bound.
Enlarge \(k_0\) to clear every denominator, in particular that of \(L\),
and fix finitely many sections realizing this jet surjection.

Spread the varieties, maps, divisors, flag, points, covering family,
and these finitely many sections over an integral finitely generated
\(\Z\)-algebra of characteristic zero. Shrink its spectrum so that the
relevant fibers and flag are smooth, the maps defining the covering
family remain dominant, and the fixed jet evaluation remains
surjective. Spread also the Harder--Narasimhan filtration of
\(\Omega_W^1|_C\). Its quotients are vector bundles on the curve; their
semistability is open, and their degrees are constant in the family.
After another shrinking, Equation~\eqref{proj:fr:generic-nef} holds on each
reduction. Such an open set has closed points in arbitrarily large
prime characteristics. We work over algebraic closures of their
residue fields, retaining the notation \(W,C,h,x,\gamma\).

In particular, all numbers in
Equations~\eqref{proj:fr:intersection-bounds} and~\eqref{proj:fr:curve-test}
are unchanged. We use the spread covering family to test effective
divisors after reduction. We do not assert that the characteristic-zero
pseudo-effective cones, or the scaling models, specialize.

Choose a fixed sufficiently ample integral divisor \(T_0\) on \(W\),
also spread with suitable sections, and for a prime \(p\) put
\begin{equation}
\label{proj:fr:bp}
 N_p=\left\lfloor\frac p{k_0}\right\rfloor,\qquad
 B_p=k_0N_pL+T_0.
\end{equation}
The difference \(B_p-pL\) belongs to a fixed finite set of rational
divisor classes. The system
\[
 |pq\xi+B_{p,1}+B_{p,2}|
\]
contains products of \(N_p\) members of \(|k_0M|\), multiplied by a
filler nonvanishing at \(z_*\). To see that one \(T_0\) suffices, write
\(p=k_0N_p+d\), where \(0\leq d<k_0\). The remainder systems have fiber
degrees \(dq\); their finitely many monomial weights require only the
finitely many line bundles \(T_0+a\,mK\) for
\(0\leq a\leq(k_0-1)q\) to have sections nonvanishing at the selected
base points. A sufficiently ample \(T_0\) has this property, and the
chosen fillers can be spread at the same time.

Products of the fixed jet sections generate jets of order
\((2n+2)k_0N_p\). Indeed each monomial of that degree or less is a
product of \(N_p\) monomials of degree at most \((2n+2)k_0\);
multiply sections with these leading monomials and then eliminate
successive higher-order terms. No factorial is divided out. Since
\[
 (2n+2)k_0N_p>(2n+1)(p-1)
\]
for all sufficiently large \(p\), the same system surjects onto the
quotient of the local ring at \(z_*\) by the \(p\)-th powers of its
\(2n+1\) regular parameters.

\subsection{Full rank away from the diagonal}

Let \(F:W\to W'\) be relative Frobenius, and put
\[
 S_0=mK,\qquad
 \mathcal E=F_*\OO_W(B_p),\qquad
 s=\rk\mathcal E=p^n,\qquad
 U_a=H^0(W,\OO_W(B_p+paS_0))
 \quad(0\leq a\leq q).
\]
Primes on line bundles indicate the base twist, so
\(F^*S'_0=pS_0\). Projection formula gives an evaluation map
\begin{equation}
\label{proj:fr:evaluation}
 \bigoplus_{a+b=q}
 U_a\otimes U_b\otimes
 \OO_{W'\times W'}(-aS'_{0,1}-bS'_{0,2})
 \longrightarrow\mathcal E\boxtimes\mathcal E.
\end{equation}
Its generic rank is \(s^2\).

Here is a coordinate verification of this assertion. Choose local
frames of the projective-bundle axes and let \(z\) be a torus coordinate
at \(z_*\), with value \(z_0\ne0\). On its Frobenius fiber the relation
is \(z^p=z_0^p\). Over the two base Frobenius fibers, the fiber-coordinate
part of the quotient is free with basis \(1,z,\ldots,z^{p-1}\).
Project the jet surjection just obtained to its coefficient of \(1\).
The monomial \(z^j\) survives precisely when \(p\) divides \(j\), in
which case it becomes the nonzero scalar \(z_0^j\).
The global section formula for the projective bundle has weights
\(0\leq j\leq pq\). Its surviving terms \(j=pa\) are exactly
\[
 H^0(W,B_p+paS_0)\otimes
 H^0(W,B_p+p(q-a)S_0).
\]
They therefore span the tensor product of the two base Frobenius
fibers. These have dimensions \(s\) each, proving the rank assertion.
The two base points used for this argument need not coincide with \(x\).

\subsection{A small rank on the diagonal}

Let
\[
 \Delta_F=W\times_{W'}W=(F\times F)^{-1}(\Delta_{W'}).
\]
On this scheme the two copies of \(pS_0\) are canonically identified,
since both descend from \(S'_0\) on the diagonal. All source twists
in Equation~\eqref{proj:fr:evaluation} restrict on \(\Delta_{W'}\) to
\(\OO_{W'}(-qS'_0)\), with these identifications pulled back to
\(\Delta_F\). Consequently every column, after restriction to
the diagonal and this common one-dimensional twist, is evaluated from a
global section on \(\Delta_F\) of
\[
 \mathcal L_p=
 \OO_{\Delta_F}(B_{p,1}+B_{p,2}+pqS_{0,1}).
\]
The rank on the diagonal is therefore at most
\(h^0(\Delta_F,\mathcal L_p)\).

The first projection \(\Delta_F\to W\) is finite flat of degree \(s\).
Filter its direct image by powers of the ideal of the reduced
diagonal. \'{E}tale-locally in smooth coordinates its algebra is
\[
 \OO_W[\delta_1,\ldots,\delta_n]/
                  (\delta_1^p,\ldots,\delta_n^p).
\]
Its graded pieces, after the line twist, are thus the vector bundles
\begin{equation}
\label{proj:fr:graded-pieces}
 \mathcal G_j=
 \OO_W(2B_p+pqS_0)\otimes A_j(\Omega_W^1),
 \qquad 0\leq j\leq u=n(p-1),
\end{equation}
where \(A_j(V)\) is the degree-\(j\) piece of the symmetric algebra of
\(V\) modulo \(p\)-th powers. This coordinate calculation is also the
canonical Frobenius filtration
\cite[Theorem~3.7]{Proj-Sun2008}. If \(e_j=\rk A_j(\Omega_W^1)\), then
\begin{equation}
\label{proj:fr:ranks}
 \sum_{j=0}^u e_j=p^n=s.
\end{equation}
The sum is \(s\), rather than \(s^2\); the second factor of \(s\) will
come from the bound for sections.

\begin{lemma}[Uniform slope bound]
\label{proj:fr:slope}
With all characteristic-zero choices fixed, uniformly for
\(0\leq j\leq n(p-1)\),
\[
 \mu_{\max}(\mathcal G_j|_C)
     \leq 7p\,l^{n-1}H^n+O(1).
\]
The constant implicit in \(O(1)\) is independent of \(p\) and \(j\).
\end{lemma}

\begin{proof}
Put \(V=\Omega_W^1|_C\). For a vector bundle on the smooth curve \(C\),
write
\[
 L_{\min}(V)=
 \lim_{e\to\infty}
 \frac{\mu_{\min}((F_C^e)^*V)}{p^e}.
\]
Langer's Frobenius-instability estimate
\cite[Corollary~2.5]{Proj-Langer2004} states that
\[
 \mu_{\min}(V)-L_{\min}(V)
       \leq \frac{(n-1)\deg A_C}{p-1}
\]
if \(A_C\) is nef and \(T_C(A_C)\) is globally generated.
The genus is fixed, so \(A_C\) can be chosen with degree bounded
independently of \(p\). Equation~\eqref{proj:fr:generic-nef} yields
\begin{equation}
\label{proj:fr:asymptotic-slope}
 L_{\min}(V)\geq-\frac{c}{p-1}
\end{equation}
with one fixed \(c\).

We need an estimate for tensor powers whose exponent grows with \(p\).
For every \(i\geq0\),
\begin{equation}
\label{proj:fr:tensor-slope}
 \mu_{\min}(V^{\otimes i})\geq iL_{\min}(V).
\end{equation}
To prove this directly, twist \((F_C^e)^*V\) by a line bundle of
degree at most
\(-\mu_{\min}((F_C^e)^*V)+2g(C)+2\) so that it is globally generated.
This follows from Serre duality and the slope criterion for the
vanishing of \(H^1\) after subtracting any point. The \(i\)-th tensor
power is then globally generated with the corresponding \(i\)-fold
twist. Pull back any quotient of \(V^{\otimes i}\), take degrees,
divide by \(p^e\), and let \(e\) tend to infinity. The bounded genus
term disappears, proving Equation~\eqref{proj:fr:tensor-slope}.
Thus every quotient of a tensor power with \(i\leq n(p-1)\) has
minimum slope at least \(-nc\).

Multiplication to the socle of the truncated polynomial algebra is
a perfect pairing in complementary degrees. The top monomial
transforms by \((\det V)^{p-1}\), so, equivariantly,
\[
 A_j(V)^*\otimes(\det V)^{p-1}\simeq A_{u-j}(V).
\]
The right side is a quotient of \(V^{\otimes(u-j)}\). Therefore
\[
 \mu_{\max}(A_j(V))\leq(p-1)\deg K_W|_C+nc.
\]
By Equations~\eqref{proj:fr:bp} and~\eqref{proj:fr:intersection-bounds},
\begin{align*}
 (2B_p+pqS_0)\cdot C
   &\leq4p\,l^{n-1}H^n+O(1),\\
 (p-1)K_W\cdot C
   &\leq\frac{2p}{r}\,l^{n-1}H^n.
\end{align*}
Since \(r>1\), their sum is bounded by the claimed expression,
with slack in its coefficient \(7\).
\end{proof}

\begin{lemma}[Uniform section bound]
\label{proj:fr:sections}
For the bundles in Equation~\eqref{proj:fr:graded-pieces},
\[
 h^0(W,\mathcal G_j)
      \leq e_jp^n\bigl(7^nH^n+o(1)\bigr),
\]
uniformly in \(j\). Consequently, for all sufficiently large \(p\),
\begin{equation}
\label{proj:fr:diagonal-rank}
 h^0(\Delta_F,\mathcal L_p)\leq\frac{s^2}{4}.
\end{equation}
\end{lemma}

\begin{proof}
Choose \(M_p=7p\,l^{n-1}H^n+c_0\), where \(c_0\) bounds the error in
Lemma~\ref{proj:fr:slope}, and put \(d_C=h\cdot C=l^nH^n\).
A bundle of rank \(e_j\) and maximal slope at most \(M_p\) has at
most \(e_j(M_p+1)\) sections: evaluation at
\(\lfloor M_p\rfloor+1\) distinct points is injective.
Its twist by \(-kh\) has no sections if \(kd_C>M_p\).

The divisor sequences on the fixed flag give
\[
 h^0(W,\mathcal G_j)
 \leq
 \sum_{k_1,\ldots,k_{n-1}\geq0}
 h^0\bigl(C,\mathcal G_j|_C
                    (-(k_1+\cdots+k_{n-1})h)\bigr).
\]
One obtains this by successively summing the restriction inequalities;
the remainders vanish for sufficiently negative ample twists at each
stage. At most
\((1+\lfloor M_p/d_C\rfloor)^{n-1}\) tuples contribute. Hence
\[
 h^0(W,\mathcal G_j)
 \leq e_j(M_p+1)(1+M_p/d_C)^{n-1}
 =e_jp^n\bigl(7^nH^n+o(1)\bigr).
\]
The powers of \(l\) cancel in the leading coefficient.
All error constants are fixed before \(p\), and the same bound applies
to every \(j\). Summing over the filtration and using
Equation~\eqref{proj:fr:ranks} gives
\[
 h^0(\Delta_F,\mathcal L_p)
 \leq s^2\bigl(7^nH^n+o(1)\bigr).
\]
Now \(7^nH^n\leq(14\epsilon)^n<1/4\) by
Equations~\eqref{proj:fr:epsilon} and~\eqref{proj:fr:intersection-bounds}.
This proves Equation~\eqref{proj:fr:diagonal-rank}.
\end{proof}

\subsection{The determinant contradiction}

Put \(R=s^2\). Choose \(R\) columns of
Equation~\eqref{proj:fr:evaluation} with nonzero determinant. Their
determinant is a nonzero section \(\sigma\) of an external product
\(\mathcal Q_1\boxtimes\mathcal Q_2\) of line bundles on \(W'\times W'\).
Indeed
\(\det(\mathcal E\boxtimes\mathcal E)
 =(\det\mathcal E)^s\boxtimes(\det\mathcal E)^s\),
and the column twists add the sums of their respective weights.
Consequently, under numerical identification with the base twists,
\begin{equation}
\label{proj:fr:slot-class}
 \frac{c_1(\mathcal Q_i)}{R}
   =\frac{c_1(\mathcal E)}s+\lambda_iS_0
   =\frac{B_p}{p}+\frac{p-1}{2p}K_W+\lambda_i mK,
 \qquad 0\leq\lambda_i\leq q.
\end{equation}
The last equality is degree-one Grothendieck--Riemann--Roch for
Frobenius \cite[Lemma~4.2]{Proj-Sun2008}:
\[
 c_1(F_*\OO_W(B_p))
   =p^{n-1}B'_p+\frac{p^n-p^{n-1}}2K_{W'}.
\]
Here numerical identification through the base-field twist must not
be confused with Frobenius pullback, which multiplies divisor classes
by \(p\).

For \(0\leq\lambda\leq q\), write
\(Q_\lambda=L+K_W/2+\lambda mK\). Direct calculation gives
\begin{equation}
\label{proj:fr:slot-domination}
 D^+-Q_\lambda
   =(r-\tfrac12-\lambda m)K+3rtA+\tfrac32E.
\end{equation}
This is pseudo-effective in characteristic zero because
\(r>qm+1\), \(K\) is pseudo-effective, and \(A,E\) are effective up
to the indicated equivalence. Pairing with the fixed class \(\gamma\)
therefore bounds \(Q_\lambda\cdot\gamma\) by
\((\pi_x^*D^+)\cdot\gamma\). This numerical inequality survives the
spread: it concerns only intersections of fixed divisors with the
fixed family. The difference between
Equation~\eqref{proj:fr:slot-class} and \(Q_{\lambda_i}\) is \(O(1/p)\)
in a fixed finite-dimensional space, uniformly in the chosen columns.

Let \(x'\in W'\) be the twist of \(x\). For a nonzero section of
\(\mathcal Q_i\), its effective divisor, strictly transformed on the
blowup at \(x'\), has nonnegative intersection with the spread moving
curve. Dividing this inequality by \(R\,J\cdot\gamma\) and using
Equations~\eqref{proj:fr:curve-test}--\eqref{proj:fr:slot-domination} yields
\begin{equation}
\label{proj:fr:slot-order}
 \frac{\ord_{x'}(\tau)}R\leq40\epsilon+o(1)
 \quad\text{for every }0\ne\tau\in H^0(W',\mathcal Q_i).
\end{equation}
This bounds all sections on the reduction, not merely the sections
spread from characteristic zero.

An external-product section can have order at \((x',x')\) at most
the sum of the two maximal slot orders. To verify this despite
possible cancellation, choose bases in each section space adapted to
its filtration by order at \(x'\). Their leading homogeneous terms
are linearly independent within each degree. Tensor products of
these terms are independent in each bidegree in the two disjoint
sets of local parameters. Thus the first nonzero homogeneous part
of a nonzero tensor cannot be cancelled beyond the sum of the slot
maxima. K\"unneth's formula identifies the external-product section
space with this tensor product. In particular,
\begin{equation}
\label{proj:fr:det-upper}
 \ord_{(x',x')}(\sigma)\leq R(80\epsilon+o(1)).
\end{equation}

On the other hand, Equation~\eqref{proj:fr:diagonal-rank} says that the
matrix of the chosen columns has rank at most \(R/4\) at \((x',x')\).
Trivialize its line bundles locally and perform invertible row and
column operations on the constant matrix. At least \(3R/4\) of the
resulting rows have every entry in the maximal ideal. Every term
of the determinant consequently has order at least \(3R/4\), so
\begin{equation}
\label{proj:fr:det-lower}
 \ord_{(x',x')}(\sigma)\geq\frac{3R}{4}.
\end{equation}
This is a pointwise maximal-ideal estimate; no stronger assertion
about an order along the whole diagonal is required.
Equations~\eqref{proj:fr:det-upper} and~\eqref{proj:fr:det-lower} contradict
\(80\epsilon<3/4\) for all sufficiently large \(p\).

We have made the choices in the order
\[
 n,\ \epsilon,\ q,\ (t,r),\
 (W,H,h,\text{flag},x,\gamma),\
 (k_0,T_0,\text{fixed sections}),\ p.
\]
In particular every constant suppressed in the estimates is fixed
before the last prime tends to infinity. This contradiction excludes
the assumed counterexample and proves
Theorem~\ref{proj:thm:smooth-nonvanishing}.
